% --- LaTeX Packages ---
% --- LaTeX Packages ---
\usepackage{amsmath}
\usepackage{amsthm}
\usepackage{bm}            % For bold math symbols
\usepackage{graphicx}      % For including images
\usepackage{xcolor}        % Required for \highlight
\usepackage{xspace}        % Required for commands like \nref
\usepackage{tikz}
%\usetikzlibrary{arrows.meta}
%\usepackage{mathtools}
%\usepackage{filecontents}
\usepackage{booktabs}
%\usepackage{placeins}
\usepackage{hyperref}
        \hypersetup{
          colorlinks=true,
          linkcolor=blue,
          urlcolor=blue,
          citecolor=blue
        }
\usepackage{etoolbox}

\setlength{\bibsep}{5pt plus 5pt} % space BETWEEN items
% For tighter lines WITHIN each item:
\usepackage{setspace}
\AtBeginEnvironment{thebibliography}{\setstretch{1.05}}

% --- Custom Commands ---

% Acronyms and Custom Text
\newcommand{\CVIC}{\text{CVIC}}
\newcommand{\deviance}{\text{dev}}
\newcommand{\dpois}{\text{dpois}}
\newcommand{\ic}{\text{ic}}
\newcommand{\nmsp}{\text{\scriptsize NMSP}}
\newcommand{\nrsp}{\text{\scriptsize NRSP}}
\newcommand{\pmin}{p_{\text{\scriptsize min}}}
\newcommand{\pp}{\text{pp}}
\newcommand{\nref}{[\textbf{references}]}
\newcommand{\SP}{survival probability\xspace}
\newcommand{\SPs}{survival probabilities\xspace}
\newcommand{\svf}{S_{i}(\cdot)}
\newcommand{\unif}[0]{\text{Unif}}
\newcommand{\rsp}[0]{\text{RSP}}
\newcommand{\msp}[0]{\text{MSP}}
\newcommand{\PPP}[0]{\text{PPP}}
\newcommand{\PPPs}[0]{\text{PPPs}}
\newcommand{\PPC}[0]{\text{PPC}}
\newcommand{\bU}[0]{\bm{U}}
\newcommand{\pv}[0]{\mathcal{P}}

% Math Symbols (Bold)
\newcommand{\btheta}[0]{\bm{\theta}}
\newcommand{\bTheta}[0]{\bm{\Theta}}
\newcommand{\bx}[0]{\bm{x}}
\newcommand{\by}[0]{\bm{y}}
\newcommand{\bY}[0]{\bm{Y}}
\newcommand{\bbeta}{\bm{\beta}}
\newcommand{\bBeta}{\bm{\Beta}}
\newcommand{\blambda}{\bm{\lambda}}
\newcommand{\bmu}{\bm{\mu}}
\newcommand{\bphi}{\bm{\phi}}
\newcommand{\bp}{\bm{p}}
\newcommand{\bsgm}{\bm{\sigma}}
\newcommand{\hu}[0]{\pi^o_{i}}
\newcommand{\mhu}[0]{\pi^o}
% General Superscripts and Subscripts
\newcommand{\obs}[0]{^{\text{obs}}}
\newcommand{\sobs}[0]{^{\text{obs}}}
\newcommand{\post}[0]{^{\text{Post}}}
\newcommand{\cv}[0]{^{\text{CV}}}
\newcommand{\iscv}[0]{^{\text{ISCV}}}
\newcommand{\supt}[0]{^{(t)}}
\newcommand{\mci}{^{(t)}}        % Kept as it has a different name from \supt
\newcommand{\mcs}{^{(s)}}
\newcommand{\rep}{^{\text{rep}}}
\newcommand{\ppc}{_{\text{ppc}}}
\newcommand{\noi}{_{-i}}
\newcommand{\wi}{_{1:n}}
\newcommand{\IS}{^{\text{\scriptsize IS}}}

% Compound Superscripts for Z-residuals
\newcommand{\rspp}[0]{^{\text{Post-RSP}}}
\newcommand{\mspp}[0]{^{\text{Post-MSP}}}
\newcommand{\rspis}[0]{^{\text{ISCV-RSP}}}
\newcommand{\mspis}[0]{^{\text{ISCV-MSP}}}
\newcommand{\rsppost}{^{\text{Post-RSP }}}  % From first block
\newcommand{\rspiscv}{^{\text{ISCV-RSP}}}    % From first block
\newcommand{\E}{\mathbb{E}}
\newcommand{\Epost}{\mathbb{E}_{\btheta\,|\,\by}}
\newcommand{\Ecv}{\mathbb{E}_{\btheta\,|\,\by_{-i}}}
\newcommand{\highlight}[1]{\textcolor{blue}{#1}}
\newcommand{\pvalue}[0]{$p$-value\xspace}
\newcommand{\pvalues}[0]{$p$-values\xspace}
\newcommand{\norm}[1]{\left\lVert#1\right\rVert}

\newcommand{\textred}[1]{\textcolor{orange}{#1}}


% ===================================================================
%                 THEOREM AND PROOF DEFINITIONS
% ===================================================================



%old theorem style for ba

% % --- 1. Define Custom Theorem Styles ---
% % A style for theorems where the optional note is bold
% \newtheoremstyle{boldnote}
%   {}		    % Space above
%   {}		    % Space below
%   {\itshape}	% Body font
%   {}		    % Indent amount
%   {\bfseries}	% Theorem head font
%   {.}		    % Punctuation after head
%   {.5em}	    % Space after head
%   {\thmname{#1}\thmnumber{ #2}\thmnote{ [\textbf{#3}]}}


% % --- 2. Define Theorem-Like Environments ---
% % Set the default style for all environments defined from this point
% \theoremstyle{plain}

% % The main theorem environment, numbered per section. All others will follow its counter.
% \newtheorem{theorem}{Theorem}[section] 

% % Environments that share the theorem counter
% \newtheorem{lemma}[theorem]{Lemma}
% \newtheorem{corollary}[theorem]{Corollary}

% % Switch to the custom style for any environments you want to use it
% % \theoremstyle{boldnote}
% % \newtheorem{specialtheorem}[theorem]{Special Theorem} % Example


% % --- 3. Define Proof and Step Environments ---
% % For custom "Step" counter within proofs
% \newtheorem{proofpart}{Step}

% % Change the name "Proof" to "Proof" (bold)
% \renewcommand{\proofname}{\textbf{Proof}}

% % Automatically reset the 'Step' counter at the beginning of every proof
% \AtBeginEnvironment{proof}{\setcounter{proofpart}{0}}
% % ===================================================================
